\thispagestyle{empty}
\addcontentsline{toc}{section}{KI-Verzeichnis}
\section*{Erklärung und Nutzungsdokumentation zum Einsatz von KI-basierten Werkzeugen bei der Anfertigung von wissenschaftlichen Arbeiten als Prüfungsleistungen}


Zur Verwendung KI-gestützter Werkzeuge erkläre ich in Kenntnis des Hinweisblatts "Hinweise zum Einsatz von KI-basierten Werkzeugen bei der Anfertigung wissenschaftlicher Arbeiten, u.a. im prüfungsrechtlichen Kontext" \space Folgendes:

\begin{itemize}[label=\textbullet, itemsep=1em] % itemsep fügt Abstand zwischen den Punkten hinzu
    \item Ich habe mich aktiv über die Leistungsfähigkeit und Beschränkungen der in meiner Arbeit eingesetzten KI-Werkzeuge informiert.
    \item Bei der Anfertigung der Arbeit habe ich durchgehend eigenständig und beim Einsatz KI-gestützter Werkzeuge maßgeblich steuernd gearbeitet.
    \item Insbesondere habe ich die Inhalte entweder aus wissenschaftlichen oder anderen zugelassenen Quellen entnommen und diese gekennzeichnet oder diese unter Anwendung wissenschaftlicher Methoden selbst entwickelt.
    \item Mir ist bewusst, dass ich als Autor/in der Arbeit die volle Verantwortung für die in ihr gemachten Angaben und Aussagen trage.
    \item Soweit ich KI-gestützte Werkzeuge zur Erstellung der Arbeit eingesetzt habe, sind diese jeweils mit dem Produktnamen, den formulierten Eingaben (Prompts), der Einsatzform sowie der entsprechenden Seiten-/Bereichsreferenzierung auf die Arbeit im KI-Verzeichnis am Ende der Arbeit vollständig ausgewiesen und im Text belegt (z.B. als Fußnote).
\end{itemize}

\vspace{2cm} % Abstand für die Unterschriftszeilen

\noindent
\begin{tabular}{p{7cm} p{7cm}}
    \rule{\linewidth}{0.4pt} & \rule{\linewidth}{0.4pt} \\
    Ort, Datum & Unterschrift der/des Studierenden Max Müller \\
\end{tabular} 
\newline
\newline
\noindent
\begin{tabular}{p{7cm} p{7cm}}
    \rule{\linewidth}{0.4pt} & \rule{\linewidth}{0.4pt} \\
    Ort, Datum & Unterschrift der/des Studierenden Jeremy Barenkamp\\
\end{tabular}
\newline
\newline
\noindent
\begin{tabular}{p{7cm} p{7cm}}
    \rule{\linewidth}{0.4pt} & \rule{\linewidth}{0.4pt} \\
    Ort, Datum & Unterschrift der/des Studierenden Anton Geiger\\
\end{tabular}
\newpage

\thispagestyle{empty}
\section*{KI-Verzeichnis-Tabelle}
Das KI-Verzeichnis ist in der vorgegebenen tabellarischen Form zu führen und abzugeben. Es empfiehlt sich, das Verzeichnis bereits ab Beginn des Erarbeitens und Schreibens der Arbeit fortlaufend zu führen. 

\begin{longtable}{| p{3.5cm} | p{4cm} | p{4cm} | p{3cm} |}
\hline
\textbf{KI-gestütztes Werkzeug} & \textbf{Beispiel-Prompts} & \textbf{Einsatzform} & \textbf{Seiten-/Bereichsangabe}  \\
\hline
\endhead % End of table header
%
% Beispiel-Einträge
Gemini 2.5 Flash & Ich habe folgenden Text geschrieben .... . Kannst du den Text pr"aziser und wissenschaftlich und grammatikalisch korrekt formulieren? & Umformulierung in präzisere Sprache und Korrektur von Rechtschreibfehler & Kapitel 1 - 9 \\
\hline
Gemini 2.5 Flash & Ich habe folgendes Tabelle für eine Planung. Bitte teile diese in 3 Tabellen auf mit jeweils einen kompletten Planungsjahr   & Generierung von Latex Code für Tabellenformatierung \\
\hline
Gemini 2.5 Flash & Ich habe folgendes Bild eines KI Verzeichnisses von meiner Hochschule, welches ich als Latex Code f"ur meine Seminararbeit brauche ... & Generierung von Latex Code f"ur dieses KI Verzeichnis & Diese und vorherige Seite \\
\hline

\end{longtable}